\pagestyle{plain}

\section{INTRODUÇÃO}
Quando olhamos para os dados sobre a maternidade no Brasil, um ponto que logo chama atenção é o crescimento contínuo das chamadas “mães solo”. Esse grupo é definido por \textcite[2]{souza2024} como "famílias monoparentais chefiadas por mulheres, que não vivem com o pai da criança e que assumem sozinhas a responsabilidade de cuidar dos filhos". 

Segundo levantamento de \textcite{feijo2023}, “entre os anos de 2012 e 2022, o número de domicílios com mães solo [no Brasil] cresceu 17,8\%, passando de 9,6 milhões para 11,3 milhões”. No Piauí, vemos essa mesma tendência se repetir. Embora dados de registro civil sem a presença do nome do pai não sejam sinônimo de abandono parental, os registros do módulo Pais Ausentes, mantido pelo Portal da Transparência do Registro Civil \cite{arpen2026}, apontam que, no Piauí, 7,5\% dos nascimentos em 2024 e 7,8\% em 2025 tiveram apenas o nome da mãe na certidão, o que pode servir de indicativo para a dimensão da situação no contexto do Estado.

Diante de um quadro como esse, torna-se necessário pensar em alternativas práticas que ajudem, ao menos em parte, a diminuir os problemas causados pela falta de uma rede de apoio. Claro que uma questão com raízes tão profundas não se resolve de forma isolada, até porque, como também destacam \textcite[1]{souza2024}:
\begin{quote}
“os resultados indicam que a maternidade solo está relacionada a desigualdades de natureza racial, educacional e territorial, e identificam-se contrastes importantes entre contextos urbanos e rurais, bem como entre diferentes níveis de instrução”.
\end{quote}

Levando em conta essas dificuldades, chegamos à seguinte questão norteadora para o trabalho: como desenvolver um aplicativo móvel seguro, baseado em Arquitetura Hexagonal e privacidade de dados, que facilite a doação de itens infantis e o apoio mútuo entre mães solo?
